\documentclass[14pt]{extreport}
\usepackage{indentfirst}
\usepackage[T2A]{fontenc}
\usepackage[utf8]{inputenc}
\usepackage[english,russian]{babel}
\usepackage{parskip}
\usepackage[top = 20mm, left = 20mm, right = 20mm, bottom = 20mm]{geometry}
\setlength\parindent{12.5mm} 
\renewcommand{\baselinestretch}{1.5} 
\def\hrf#1{\hbox to#1{\hrulefill}} 
\usepackage{ragged2e} 
\usepackage{caption}
\usepackage{tikz}
\usepackage{circuitikz}
\usetikzlibrary{decorations.pathmorphing,arrows.meta}
\usepackage{multicol}
\usepackage{float}
\usepackage{tipa}
\usepackage{graphicx}
\usepackage{enumitem}
\usepackage{amsmath}
\usepackage{amssymb}
\usepackage{listings}
\usepackage{moreverb}
\usepackage{minted}
\setcounter{tocdepth}{3}
\setcounter{secnumdepth}{3}

\lstset{
   breaklines=true,
   basicstyle=\ttfamily,
   extendedchars=\true,
   keepspaces=true,
   extendedchars=\true}
\usepackage[colorlinks=true, citecolor=blue, urlcolor=blue, linkcolor=blue]{hyperref}
\renewcommand\thesection{\arabic{section}}


\begin{document}

\thispagestyle{empty}
    \centering{
    \hspace{0.2cm}МИНИСТЕРСТВО НАУКИ И ВЫСШЕГО ОБРАЗОВАНИЯ РОССИЙСКОЙ ФЕДЕРАЦИИ\\
    Федеральное государственное автономное образовательное учреждение высшего образования\\
    {\bf «Санкт-Петербургский политехнический университет Петра Великого»}\\[10pt]
    Институт компьютерных наук и кибербезопасности\\
    Высшая школа технологий искусственного интелекта\\
    02.03.01 Математика и компьютерные науки\\
    \vspace{1.5cm}
    {\large
    \bf Отчёт\\
    \bf Современные технологии разработки программного обеспечения}\\
    {\large\bf Отладка и анализ ошибок в C++}\\
    \vspace{\fill}
    }
    \vspace{1.5cm}
    \flushleft{
        {\bf Выполнил:}\\
        студент группы 5130201/50301 \hspace{2.25cm} 
        \line(1,0){120} \quad Ташлыкова Е.В.\\
        {\bf Принял:}\\
        Ассистент \hspace{66mm}
        \line(1,0){120} \quad Кирпиченко С.Р.\\
        \vspace{1cm}
        \hspace{11.6cm}<<\hrf{2em}>> \hrf{6em} 2026г.
        }
        
    \vspace{1.5cm}
    \centering{
    Санкт-Петербург, 2026
    }

\newpage
\justifying
{\centering\section*{Аннотация}}
Данная работа посвящена исследованию типовых ошибок в языке программирования C++ и методов их обнаружения с использованием современных инструментов отладки и анализа. 

\textbf{Цель работы} — систематизация знаний о причинах возникновения типовых ошибок в С++ и анализ различных инструментов их обнаружения. 

\textbf{Методы:} воспроизведение минимальных примеров ошибочного кода, динамический анализ с использованием Valgrind, AddressSanitizer (ASan), \\
\noindent UndefinedBehaviorSanitizer (UBSan), статический анализ с помощью clang-tidy. 

\textbf{Результаты:} приведены и проанализированы причины возниконовения segmentation fault (в том числе частные случаи UB), утечки памяти, ошибки линковки, undefined behavior; приведены примеры обнаружения ошибок. 

\textbf{Область применения:} результаты могут быть использованы в учебном процессе и практической разработке на С++ для повышения надёжности ПО.

% В работе рассмотрены четыре категории ошибок: ошибка сегментации (Segmentation fault), утечка памяти, ошибка линковки и неопределённое поведение (Undefined behaviour). Для каждого типа ошибки приведены минимальные воспроизводимые примеры кода, объяснены причины возникновения и продемонстрированы способы обнаружения с помощью инструментов динамического анализа (Valgrind, AddressSanitizer, UndefinedBehaviorSanitizer) и статического анализа (clang-tidy). Особое внимание уделено обязательному использованию статического анализатора, что позволяет выявлять ошибки на этапе компиляции без выполнения программы. Результаты работы могут быть использованы для повышения качества и надёжности программного обеспечения на C++.
\newpage
\renewcommand\contentsname{\vspace{-3.2cm}\centering\LargeСодержание}
\tableofcontents
\thispagestyle{empty}
\newpage
{\centering\section*{Введение}}\addcontentsline{toc}{section}{Введение}
Язык C++ предоставляет разработчику высокий уровень контроля над памятью и аппаратными ресурсами, однако это накладывает дополнительную ответственность на порграммиста за корректное управление ресурсами. Ошибки в работе с памятью, некорректные операции и проблемы линковки могут приводить к аварийным завершениям программ, утечкам ресурсов и неопределённому поведению.

Для обнаружения подобных ошибок существуют специализированные инструменты: 
\begin{itemize}
    \item Статический анализ (clang-tidy, PVS-Studio): проверяет исходный код без выполнения. 
    \item Динамический анализ (Valgrind, AddressSanitizer, UndefinedBehaviorSanitizer): отслеживает ошибки во время работы программы. 
\end{itemize}

%\newpage
%{\centering\section{Задачи}}
%Текст

\newpage
{\centering\section{Segmentation fault}}
Ошибка, которая возникает во время обращения программы к участку памяти, доступ к которой ей запрещен или которая не существует. \\
Частые причины возникновения ошибки:
\begin{itemize}
    \item Получение доступа к памяти, которой программа не владеет.
    \item Разыменовывание nullptr указателя.
    \item Изменение памяти доступпной только для чтения.
\end{itemize}

В стандарте языка C++ отсутствует понятие «segmentation fault». Большинство действий, приводящих к этой проблеме, согласно стандарту C++, являются неопределённым поведением (Undefined Behaviour, UB). Segmentation fault -- это одно из возможных проявлений UB на платформах с менеджером памяти и защитой страниц.

\hspace{0.5cm}

\centerline{\large{\textbf{1.1 Изменение строкового литерала}}}

\begin{minted}[fontsize=\small]{c}
int main()
{
    char* str;
    str = "hello";
    *(str + 1) = 'n'; // Error
    return 0;
}
\end{minted}

Проблема возникает потому что мы пытаемся записать данные в память, доступную только для чтения. Попытка записи -- UB, на практике -- сегфолт.

\hspace{0.5cm}

\centerline{\large{\textbf{1.2 Обращение к освобожденной памяти}}}
\begin{minted}[fontsize=\small]{c}
int main()
{
    int n = 10;
    char *str = new char[n];
    str = "hello";
    delete str;
    *str = "world"; // Error
    return 0;
}
\end{minted}
Проблема возникает потому что, указатель str разыменовывается после освобождения блока памяти, что недопустимо с точки зрения компилятора. Это UB, часто -- сегфолт или повреждение данных.

\hspace{0.5cm}

\centerline{\large{\textbf{1.3 Разыменовывание нулевого указателя}}}
\begin{minted}[fontsize=\small]{c}
#include <iostream>
int main()
{
    int* ptr;
    std::cout << *ptr; // Error
    return 0;
}
\end{minted}
Происходит обращение к неизвестному участку памяти, что приводит к ошибке сегментации или к другому неопределенному поведению. 

\hspace{0.5cm}

\centerline{\large{\textbf{1.4 Выход за границы массива}}}
\begin{minted}[fontsize=\small]{c}
int main()
{
    int arr[2];
    arr[3] = 10; // Error
    return 0;
}
\end{minted}
Обращение к элементу массива по индексу, превышающему его размер, приводит к чтению или записи памяти, не принадлежащей массиву. Из-за этого могут возникнуть ошибки сегментации или другое неопределенное поведение. 

Обнаружить ошибки можно с помощью AddressSanitizer: 
\begin{verbatim}
g++ -fsanitize=address -g example.cpp -o example
./example
\end{verbatim}
ASan поймает выход за границы, use-after-free и разыменование нуля, указав точную строку.

\newpage
{\centering\section{Утечка памяти}}
Это ситуация, когда память, выделенная для выполнения определенной задачи, остается выделенной даже после того, как необходимость в ней отпадает. Это приводит к нерациональному использованию памяти, поскольку она недоступна для других задач до завершения работы программы.
Так как в С++ нет автоматических сборщиков мусора любое выделение и освобождение динамической памяти становится ответственностью программиста.

\begin{minted}[fontsize=\small]{c}
void f() 
{
    int* ptr = new int[10]; // Выделяем память
    // Не освобождаем ее 
    return;
}
int main() 
{
    f();
}
\end{minted}
Выделенная память перестает быть доступной, но при этом остается занятой. Со временем, если программа продолжает расходовать память, то система начнет работать медленнее, появятся доругие ошибки, а память, так как это ограниченный ресурс, закончится.

Для обнаружения утечек памяти существуют такие иснтрументы, как Valgrind и AddressSanitizer.

Valgrind — это фреймворк для отладки и профилирования памяти. Он может обнаруживать утечки памяти, некорректные обращения к памяти и другие ошибки. Valgrind предоставляет подробные отчеты об утечках памяти, включая информацию об их происхождении и возможных причинах.
\begin{verbatim}
# установка
brew install valgrind

# компиляция с отладочной информацией и запуск 
g++ -g -O0 main.cpp -o programm
valgrind --leak-check=full ./programm
\end{verbatim}

AddressSanitizer, также известный как (Asan), — это быстрый детектор ошибок в работе с памятью, который встроен в компиляторы GCC и Clang.
\begin{verbatim}
g++ -fsanitize=address -g -O0 main.cpp -o programm && ./programm
# или
clang++ -fsanitize=address -g -O0 main.cpp -o program ./programm
\end{verbatim}

Эти интсрументы имеют свои отличия. \\
Например, \textbf{скорость:} ASan работает значительно быстрее благодаря тому, что он внедряет проверки непосредственно в код на этапе компиляции, программа замедляется, по сравнению с обычным запуском. Valgrind же эмулирует выполнение каждой инструкции процессора, что приводит к большему замедлению. 
\textbf{Потребление памяти:} Valgrind потребляет умеренное количество дополнительной памяти, тогда как ASan требует значительно больше ресурсов, потому что создаёт "теневую память" (shadow memory) для отслеживания состояния каждого байта. 
\textbf{Необходимость перекомпиляции:} Valgrind работает с уже скомпилированным бинарным файлом, достаточно запустить программу под Valgrind. ASan же требует перекомпиляции всего проекта со специальными флагами компилятора '-fsanitize=address -g'.\\
Таким образом для CI/CD и ежедневной разработки лучше подходит ASan (быстрее). Для финального анализа сложных багов или при невозможности перекомпиляции -- Valgrind. 

\newpage
{\centering\section{Ошибка линковки}}
Связывание - это процесс объединения отдельных объектных файлов и библиотек в одну исполняемую программу. Ошибки связывания возникают, когда компоновщик не может разрешить ссылки между различными объектными файлами или библиотеками. Эти ошибки препятствуют созданию конечного исполняемого файла.

Частые причины возниконовения ошибок линковки - это отсутствие реализации функции, некорректные пути к библиотекам, несовпадение сигнатур функций и циклические зависимости.
\begin{minted}[fontsize=\small]{c}
// header.hpp
int calculate(int x);

// main.cpp
#include "header.hpp"
int main() 
{
    int result = calculate(10);  // Error
    return 0;
}
\end{minted}
Возникает ошибка - неопределенная ссылка на calculate. Компоновщик не может найти её тело.


\begin{minted}[fontsize=\small]{c}
void print_val(int x) {}
void print_val(double x); 

int main() 
{
print_val(3.14);  // Error
return 0;
}
\end{minted}

Мы хотим использовать функцию с другим типом параметра. Компоновщик ищет функцию с сигнатурой print\_val(double), но существует только print\_val(int).


Для обнаружения ошибок можно использовать флаги компилятора и компоновщика:
\begin{verbatim}
gcc -Wall -Wextra -Werror main.c -o program
\end{verbatim}

\newpage
{\centering\section{Undefined behaviour}}
Неопределённое поведение — это ситуации, при которых стандарт языка C++ не определяет ожидаемое поведение программы. Компилятор может сгенерировать любой код, программа может работать некорректно, падать или делать непредсказуемые вещи.

\begin{minted}[fontsize=\small]{c}
// 1. Использовать объект после move
std::string x = "hello";
std::string y = std::move(x);
std::cout << x;  // Error
\end{minted}

\begin{minted}[fontsize=\small]{c}
// 2. Деление на ноль
int n = 0;
return n / 0; // Error
\end{minted}

\begin{minted}[fontsize=\small]{c}
// 3. Выход за пределы индексирования
int arr[4] = {0, 1, 2, 3};
return arr[5]; // Error
\end{minted}

Проблемы можно обнаружить с помощью UndefinedBehaviorSanitizer (UBSan):
\begin{verbatim}
g++ -fsanitize=undefined -g example.cpp -o example
./example
\end{verbatim}
Однако, UBSan может не поймать ошибку при использовании объекта после move, так как формально чтение перемещённого объекта не всегда является UB (зависит от типа). Для таких случаев рекомендуется статический анализ. 
Например, с помошью clang-tidy:
\begin{verbatim}
clang-tidy example.cpp --checks='*' --
\end{verbatim}

\newpage
\section*{Заключение}\addcontentsline{toc}{section}{Заключение}
В ходе выполнения работы были рассмотрены четыре основных типа ошибок в языке C++: ошибка сегментации (Segmentation fault), утечка памяти, ошибка линковки и неопределённое поведение (Undefined behaviour). Для каждого типа были приведены минимальные примеры кода, воспроизводящие соответствующую проблему, и продемонстрированы способы их обнаружения с использованием современных инструментов анализа. 

Основные выводы: 

1. \textbf{Segmentation fault} возникает при попытке доступа к памяти, которая не принадлежит программе или доступна только для чтения. AddressSanitizer эффективно обнаруживает такие ошибки, указывая точное место нарушения.

2. \textbf{Утечка памяти} происходит, когда динамически выделенная память не освобождается после её использования. Valgrind позволяет точно идентифицировать утечки с указанием места выделения памяти.

3. \textbf{Ошибка линковки} проявляется на этапе сборки программы и связана с неразрешёнными ссылками между объектными файлами. Флаги компилятора `-Wall -Wextra -Werror` помогают выявить многие проблемы на ранних стадиях.

4. \textbf{Неопределённое поведение} — наиболее опасный тип ошибок, так как программа может работать некорректно без видимых сообщений. \\
\noindent UndefinedBehaviorSanitizer позволяет выявить многие случаи UB, включая деление на ноль и использование перемещённых объектов.

5. \textbf{Статический анализ} (clang-tidy) является обязательным инструментом в современной разработке, так как позволяет находить ошибки без выполнения программы, что особенно необходимо для крупных проектов.

% Практическая значимость работы заключается в формировании подхода к отладке C++-программ, который сочетает статический и динамический анализ для максимального покрытия возможных ошибок. Рекомендуется включать эти инструменты в процесс непрерывной интеграции (CI/CD) для автоматической проверки качества кода.

\newpage
\renewcommand{\bibname}{\vspace{-3.2cm}\centering\LargeСписок используемых источников}\addcontentsline{toc}{section}{Список используемых источников}
\begin{thebibliography}{8}

\bibitem{dyuhurst}
Дьюхэрст С. Скользкие места C++. Как избежать проблем при проектировании и компиляции ваших программ / С. Дьюхэрст. — Москва : ДМК Пресс, 2006. — 264 с. — ISBN 5-94074-083-9.

\bibitem{geeks_segfault}
GeeksforGeeks. "Segmentation Fault in C/C++" [Электронный ресурс]: 
\href{https://www.geeksforgeeks.org/segmentation-fault-c-cpp/}{https://www.geeksforgeeks.org/segmentation-fault-c-cpp/} 
(дата обращения: 10.05.2026).

\bibitem{geeks_memory_leak}
GeeksforGeeks. "Memory Leak in C++ and How to Avoid It" [Электронный ресурс]: 
\href{https://www.geeksforgeeks.org/memory-leak-in-c-and-how-to-avoid-it/}{https://www.geeksforgeeks.org/memory-leak-in-c-and-how-to-avoid-it/} 
(дата обращения: 10.05.2026).

\bibitem{geeks_detect_memory}
GeeksforGeeks. "Detect Memory Leaks in C++" [Электронный ресурс]: 
\href{https://www.geeksforgeeks.org/detect-memory-leaks-in-cpp/}{https://www.geeksforgeeks.org/detect-memory-leaks-in-cpp/} 
(дата обращения: 11.05.2026).

\bibitem{labex_linking}
LabEx. "How to Handle Linking Errors Effectively" [Электронный ресурс]: 
\href{https://labex.io/ru/tutorials/c-how-to-handle-linking-errors-effectively-419526}{https://labex.io/ru/tutorials/c-how-to-handle-linking-errors-effectively-419526} 
(дата обращения: 11.05.2026).

\bibitem{geeks_ub}
GeeksforGeeks. "Undefined Behavior in C/C++" [Электронный ресурс]: 
\href{https://www.geeksforgeeks.org/undefined-behavior-c-cpp/}{https://www.geeksforgeeks.org/undefined-behavior-c-cpp/} 
(дата обращения: 12.05.2026).

\bibitem{clangtidy}
Clang-Tidy — инструмент статического анализа для C++ [Электронный ресурс]: \href{https://clang.llvm.org/extra/clang-tidy/}{https://clang.llvm.org/extra/clang-tidy/} 
(дата обращения: 12.05.2026).

\end{thebibliography}

\end{document}